\chapter{Formalization Map: Reusable Pro-etale and Perverse APIs}
\label{app:proetale-formalization}

This appendix is a theorem-boundary specification, not Lean code.  It records
an implementation order in which each layer is reusable outside the Weil
conjectures and in which the final theorem cannot acquire a hidden
project-specific axiom.

\section{The starting point and the first audit}

The publicly documented mathlib endpoint already contains the small
pro-etale site and an \(\ell\)-adic coefficient sheaf with its sheaf
cohomology; see
\href{https://leanprover-community.github.io/mathlib4_docs/Mathlib/AlgebraicGeometry/Sites/ElladicCohomology.html}
{mathlib's \texttt{ElladicCohomology} module}.
This is a genuine head start, but it is not yet the constructible
six-functor category used by the proof.

At the pinned project commit, the first mechanical audit should produce a
minimal-import file checking at least:

\begin{enumerate}
\item weakly etale morphisms, the pro-etale site, and its functoriality;
\item the exact type and universe of the existing coefficient sheaf;
\item the ambient sheaf category and its Grothendieck/abelian instances;
\item the available derived category, \(t\)-structure, and localization
      infrastructure;
\item the definitions of sheaf cohomology and the existing
      \(\ell\)-adic cohomology endpoint.
\end{enumerate}

No declaration should be inferred from a similarly named documentation page.
The inventory consists of successful \(\#check\) commands under the minimal
imports used by the eventual file.

\section{Layer 0: coefficient-independent categorical infrastructure}

The first layer should not mention schemes, finite fields, or weights.  It
exports:

The completion identities in this layer are parameterized by hypotheses;
they are not unconditional facts about an arbitrary ringed topos.  The
reconstruction theorem used in the proof spine assumes a replete topos, a
noetherian coefficient ring, and a finitely generated ideal.  In the more
general algebraic API, the ideal must be weakly proregular for the standard
derived-completion formulas.  Comparisons with compatible systems in an
unbounded derived category require the stronger regular-ideal hypothesis
appearing in Bhatt--Scholze, Lemma 3.5.7(5).  For the complete DVR used here,
the principal ideal \((\varpi)\) satisfies these conditions.  No global
instance should silently infer them for arbitrary rings or topoi.

\begin{enumerate}
\item derived completion for a commutative ring object \(R\) and an ideal
      generated by \(\varpi\):
      \[
       \operatorname{complete}_\varpi(K)
        =\Rlim_n(K\otimes_R^LR/\varpi^n);
      \]
\item the predicate \(\operatorname{DerivedComplete}_\varpi(K)\);
\item idempotence of completion and the universal property as a left
      adjoint to inclusion of complete objects;
\item conservativity of all finite reductions on complete objects satisfying
      a separatedness or constructibility hypothesis;
\item completed pullback and completed tensor as generic constructions;
\item dualizable objects in a closed symmetric monoidal derived category;
\item transport of dualizability and completion across equivalences.
\end{enumerate}

\begin{formalization}
Suggested public interfaces include
{\small
\[
\begin{aligned}
\texttt{derivedComplete}(K)&:\operatorname{Prop},\\
\texttt{completion}(K)&:\D(R),\\
\texttt{completionIsoRlim}(K;
  \texttt{Replete},\texttt{Noetherian},\texttt{Ideal.FG})&:
 K^\wedge\simeq\Rlim_n(K\otimes^LR/\varpi^n),\\
\texttt{unboundedSystemComparison}
  (\texttt{Ideal.IsRegular})&:\operatorname{Equivalence},\\
\texttt{isIsoOfReductionIsIso}(f)&:
 \bigl(\forall n,\operatorname{IsIso}(f_n)\bigr)
 \to\operatorname{IsIso}(f).
\end{aligned}
\]
}
A generic commutative-algebra version of the first equivalence should instead
take \texttt{Ideal.WeaklyProregular}.  The last theorem must state the
completeness and separatedness or constructibility hypotheses under which
finite reductions are conservative; it is not true for arbitrary
noncomplete objects.
\end{formalization}

\section{Layer 1: the pro-etale coefficient sheaf}

For a scheme \(X\) and complete noetherian local ring \(\mathcal O\), define
the sheaf associated with the topological ring:
\[
 \widehat{\mathcal O}_X(U)
   =C_{\mathrm{cont}}(\pi_0(U),\mathcal O)
\]
on the weakly contractible basis, followed by sheafification or the
equivalent topological constant-sheaf construction.  Export canonical
isomorphisms
\[
 \widehat{\mathcal O}_X
   \simeq\varprojlim_n(\mathcal O/\varpi^n)_X,
\qquad
 f^{-1}\widehat{\mathcal O}_Y
   \longrightarrow\widehat{\mathcal O}_X.
\tag{C.1}\label{eq:formalization-coefficient-sheaf}
\]

The existing \(\Zl\)-valued sheaf should be identified with the specialization
of this generic construction, not replaced by a second unrelated constant
sheaf.  The coefficient interface should accept \(\mathcal O_E\) as well as
\(\Zl\); otherwise the Artin--Schreier kernel will force a redesign later.

\section{Layer 2: lisse and constructible objects}

The reusable derived definition is
\[
 \operatorname{Lisse}(K)
   :\Longleftrightarrow
   \operatorname{IsDualizable}(K).
\tag{C.2}\label{eq:formalization-lisse-dualizable}
\]
On a weakly contractible affine this must be compared with a perfect complex
of coefficient modules.  Define constructibility by the existence of a
finite locally closed stratification on which all restrictions are lisse.
The restriction of the ambient standard \(t\)-structure is installed only
under the hypotheses of Hemo--Richarz--Scholbach, Theorem 6.2: the condensed
coefficient ring is \(t\)-admissible; for lisse objects, \(X\) has Zariski
locally finitely many irreducible components; and for constructible objects,
every constructible subset of \(X\) has Zariski locally finitely many
irreducible components.  The coefficient-limit and comparison interfaces
separately retain their qcqs and transition-map hypotheses.  Locally
noetherian schemes and the coefficients \(\mathcal O_E,E\) used in these
notes meet the selected conditions, but there is no such instance for an
arbitrary scheme and arbitrary condensed ring.

For the integral category, package the equivalent interface
\[
\begin{aligned}
\operatorname{IntegralConstructible}(K)
 :\Longleftrightarrow{}&
 \operatorname{DerivedComplete}_\varpi(K)\\
 &{}\land
 \operatorname{Constructible}
   (K\otimes^L\mathcal O/\varpi)\\
 &{}\land
 \operatorname{FiniteTorAmplitude}(K).
\end{aligned}
\tag{C.3}\label{eq:formalization-integral-constructible}
\]
The equivalence with stratified local perfectness is a theorem, not a second
definition exposed to downstream users.

\begin{formalization}
The layer should export:
\begin{enumerate}
\item stability under shifts, cones, retracts, tensor, and internal Hom;
\item restriction to locally closed subschemes;
\item common refinement for finitely many constructible objects;
\item geometric-stalk conservativity on constructible objects;
\item finite-level reduction and reconstruction by \(\Rlim\);
\item coefficient extension
      \(\mathcal O\to E\to E'\) and descent of finite diagrams from
      \(\Qbarl\) to finite \(E\).
\end{enumerate}
Do not store a chosen stratification in the public object.  Store its
existence and expose choice-independent restriction theorems.
\end{formalization}

\section{Layer 3: comparison with classical adic systems}

The primary API remains the pro-etale category.  The classical tower model is
a comparison target:
\[
 \texttt{adicSystemRealization}:
 \mathrm D^b_{c,\mathrm{adic}}(X,\mathcal O)
 \xrightarrow{\sim}
 \Dcons(X_{\proet},\widehat{\mathcal O}).
\tag{C.4}\label{eq:formalization-adic-comparison}
\]

Its unit and counit are the reconstruction maps
\[
 \Rlim_n\nu^*K_n\longrightarrow K,
\qquad
 \bigl(\Rlim_n\nu^*K_n\bigr)/\varpi^m\longrightarrow\nu^*K_m.
\]
All later compatibility statements should be phrased as natural
isomorphisms involving this equivalence.  The project should not implement
the AR-null quotient as the primary public category merely because it is the
historical model.

\section{Layer 4: ambient solid operations and constructible restriction}

Route B separates ambient existence from geometric preservation:

\begin{enumerate}
\item the unrestricted pro-etale category and the solidification reflector;
\item a presentable closed monoidal category \(\Dsolid(X,\Lambda)\);
\item the correspondence-level ambient operations
      \(f^*,f_*,f_!,f^!,\otimes,R\mathcal Hom\);
\item the full subcategory \(\Dindcons\) and the theorem
      \(\Dcons=\Dindcons^\omega\);
\item an intrinsic constructibility predicate, defined by finite
      stratifications and dualizability;
\item separate geometric theorems that every ambient operation preserves
      constructibility under explicit hypotheses.
\end{enumerate}

The ambient functor record must not contain constructibility hypotheses.
Those belong to restriction lemmas whose conclusions are functors between
the full constructible subcategories.

The exceptional direct image requires a choice-management structure:
\[
\begin{array}{l}
\texttt{Compactification}(f),\\
\texttt{rfShriekUsing}(c,K),\\
\texttt{rfShriekComparison}(c_1,c_2,K),\\
\texttt{rfShriekComparison\_cocycle},\\
\texttt{rfShriek}(f,K)
\quad\text{with a canonical-isomorphism specification}.
\end{array}
\tag{C.5}\label{eq:formalization-rfshriek}
\]
If a canonical choice of compactification is made for implementation, no
theorem statement may depend on that choice.

Integral constructible functors should still export finite-reduction
compatibilities.  For example,
\[
\begin{aligned}
\texttt{reduction\_rfStar}:&
 (Rf_*K)_n\simeq Rf_*(K_n),\\
\texttt{reduction\_rfShriek}:&
 (Rf_!K)_n\simeq Rf_!(K_n),\\
\texttt{reduction\_rfUpperShriek}:&
 (Rf^!L)_n\simeq Rf_n^!(L_n).
\end{aligned}
\tag{C.6}\label{eq:formalization-six-reduction}
\]
These are comparison theorems, not definitions.  Base change, projection
formula, Kunneth, localization, and their higher coherence are first fields
or theorems of the ambient correspondence formalism and then restrict to
constructibles.  The finite-reduction API remains useful as a conservative
test and as the bridge to classical torsion geometry; it is not the sole
coherence engine.

\begin{auditbox}
There must be no declaration identifying \(Rf_!\) with a derived functor of
an abelian pro-etale \(f_!\).  There must also be no theorem asserting proper
base change on the unrestricted naive pro-etale category, identifying
solidity with derived completion, or asserting
\(\Dcons=(\Dsolid)^\omega\).  Ambient compactification coherence and every
constructible-preservation theorem are separate release gates.  The
finite-coefficient theorems remain dependencies of the known proof even
though the ambient operations are not defined levelwise.
\end{auditbox}

\section{Layer 5: Weil structures, duality, and traces}

A Weil object should be a structure on a geometric-base-change object:
\[
 \texttt{WeilObject}(K)
   =\{\phi:F_X^*K\xrightarrow{\sim}K\}.
\tag{C.7}\label{eq:formalization-weil-object}
\]
Iteration must be an ordered composite, with lemmas for zero, successor, and
addition.  The six operations lift to Weil objects using explicit exchange
isomorphisms.

Keep distinct types or named constructions for:

\begin{enumerate}
\item arithmetic Galois Frobenius;
\item geometric Galois Frobenius;
\item the scheme Frobenius correspondence;
\item local closed-point Frobenius;
\item the induced endomorphism of a stalk or cohomology object.
\end{enumerate}

The trace layer should first be generic for dualizable objects in a symmetric
monoidal category, then specialized to perfect coefficient complexes.
The geometric trace formula returns equality between the categorical global
trace and local terms.  The point-counting corollary is a separate theorem
that performs the fixed-point and closed-point exponent conversion.

\section{Layer 6: perverse structures, including the integral pair}

Build the abstract machinery before mentioning sheaves:

\begin{enumerate}
\item bounded \(t\)-structures and their hearts;
\item gluing along recollement;
\item one-sided and two-sided \(t\)-exact functors;
\item truncation towers and finite filtrations;
\item images in a heart and the generic intermediate-extension construction;
\item tilting a \(t\)-structure by a torsion pair.
\end{enumerate}

For a noetherian space with dimension function \(\delta\), define the rational
support and cosupport predicates and prove that they produce the middle
perverse structure.  The terminal identities are
\[
\begin{aligned}
\texttt{perverseLE}(K)&\leftrightarrow
 \forall x\,a,\ a>-\delta(x)\to H^a(i_x^*K)=0,\\
\texttt{perverseGE}(K)&\leftrightarrow
 \forall x\,a,\ a<-\delta(x)\to H^a(i_x^!K)=0.
\end{aligned}
\tag{C.8}\label{eq:formalization-perverse-predicates}
\]

Over \(\mathcal O\), define both \(p\) and \(p^+\).  The required release
theorems include:
\[
\begin{aligned}
\texttt{verdierPToPPlus}&:
 ({}^p\Perv(X,\mathcal O))^{\op}
   \xrightarrow{\sim}{}^{p^+}\Perv(X,\mathcal O),\\
\texttt{verdierPPlusToP}&:
 ({}^{p^+}\Perv(X,\mathcal O))^{\op}
   \xrightarrow{\sim}{}^p\Perv(X,\mathcal O),\\
\texttt{rationalization\_perverse}&:
 {}^p\Perv(\mathcal O)\to\Perv(E),\\
\texttt{middleExtension\_p\_rationalize}&:
 ({}^pj_{!*}P)\otimes E\simeq j_{!*}(P\otimes E),\\
\texttt{middleExtension\_pplus\_rationalize}&:
 ({}^{p^+}j_{!*}P)\otimes E\simeq j_{!*}(P\otimes E).
\end{aligned}
\tag{C.9}\label{eq:formalization-integral-perverse}
\]
Do not install a Verdier self-duality instance for the integral \(p\)-heart.

\section{Layer 7: weights and the perverse engine}

Keep the following predicates distinct:
\[
\begin{array}{ll}
\texttt{PunctuallyPure}(\iota,w,\cF),&
\texttt{WeightsLE}(\iota,w,K),\\
\texttt{WeightsGE}(\iota,w,K),&
\texttt{Pure}(\iota,w,K),\\
\texttt{UniversallyPure}(w,K).&
\end{array}
\tag{C.10}\label{eq:formalization-weight-predicates}
\]

The ordinary/perverse bridge is a first-class theorem:
\[
\operatorname{WeightsLE}(K,w)
\Longleftrightarrow
\forall a,\ 
\operatorname{WeightsLE}(\pH^aK,w+a),
\tag{C.11}\label{eq:formalization-perverse-weight-bridge}
\]
and similarly for lower bounds and purity.  Tests must verify:
\[
 \wt(K[1])=\wt(K)+1,\qquad
 \wt(K(1))=\wt(K)-2,\qquad
 \wt(E_X[d])=d
\]
when \(X\) is smooth of pure dimension \(d\).

The main theorem should be exported first in perverse form:
\[
\texttt{perverseRfShriekWeightsLE}:
\operatorname{WeightsLE}(\cP,w)
\to
\operatorname{WeightsLE}(\pH^aRf_!\cP,w+a).
\tag{C.12}\label{eq:formalization-perverse-direct-image}
\]
Its dependencies are separate declarations for quasi-finite exactness,
open--closed devissage, relative Noether normalization, affine-line
projection, Fourier exactness, and the curve estimate.  The derived theorem
is then obtained from the finite perverse truncation tower.

\begin{powerbox}
This API order exposes the mathematical gain: the hard theorem returns
canonical perverse cohomology objects with support and boundary behavior.
Ordinary cohomology is recovered only after applying the perverse weight
criterion.
\end{powerbox}

\section{Layer 8: proper purity, IC, and decomposition}

Implement the following chain as separate theorems:
\[
\begin{aligned}
\texttt{rfStarWeightsGE}&:
 \operatorname{WeightsGE}(K,w)\to
 \operatorname{WeightsGE}(Rf_*K,w),\\
\texttt{properPushforwardPure}&:
 \operatorname{Proper}(f)\to
 \operatorname{Pure}(K,w)\to
 \operatorname{Pure}(Rf_*K,w),\\
\texttt{middleExtensionPure}&:
 \operatorname{Pure}(P,w)\to
 \operatorname{Pure}(j_{!*}P,w),\\
\texttt{intersectionComplexPure}&:
 \operatorname{PunctuallyPure}(L,u)\to
 \operatorname{Pure}(\IC_X(L),u+d),\\
\texttt{properIntersectionCohomologyPure}&:
 \operatorname{Proper}(X)\to
 \operatorname{Pure}(IH^r(X,L),u+d+r).
\end{aligned}
\tag{C.13}\label{eq:formalization-proper-purity-chain}
\]

Only after these declarations may the library add geometric semisimplicity,
splitting of a pure complex, and the decomposition theorem.  The splitting
returns existence of an isomorphism after geometric base change; it is
neither canonical nor a claim that Frobenius acts semisimply.

\section{Layer 9: final assembly}

The terminal purity term for a smooth proper \(d\)-fold should visibly
compose
\[
\begin{tikzcd}
 E_X[d]
 \arrow[r,"\text{proper purity}"]&
 Ra_*E_X[d]
 \arrow[r,"H^{i-d}"]&
 H^i(X,E).
\end{tikzcd}
\tag{C.14}\label{eq:formalization-terminal-term}
\]
The full Weil theorem combines this with three independent, reusable
geometry-free modules:

\begin{enumerate}
\item trace-to-zeta via the determinant--trace identity;
\item duality-to-functional-equation via perfect similitude pairings;
\item rational-series integrality and universal-weight separation.
\end{enumerate}

\section{Milestones and release gates}

\begin{longtable}{
 >{\raggedright\arraybackslash}p{0.11\textwidth}
 >{\raggedright\arraybackslash}p{0.49\textwidth}
 >{\raggedright\arraybackslash}p{0.30\textwidth}}
\toprule
Gate & Required artifact & Mandatory tests\\
\midrule
\endhead
M0 &
Pinned inventory and minimal imports. &
All \(\#check\) declarations compile.\\

M1 &
Generic completion and coefficient sheaf. &
\(\Zl/\ell^n\simeq\mathbf Z/\ell^n\);
completion idempotence.\\

M2 &
Constructible integral and rational categories. &
Point, finite etale scheme, and lisse representation examples.\\

M3 &
Completed four operations and localization. &
Finite reduction, composition, open--closed triangles.\\

M4 &
\(Rf_!,Rf^!\), base change, duality. &
Compactification independence and cocycle; affine-line compact support.\\

M5 &
Weil objects, twists, and trace. &
Eigenvalues \(1,q^{-1},q\) on a point, \(E(1)\), and
\(H^2(\PP^1,E)\).\\

M6 &
Rational and integral perverse structures. &
\(p/p^+\) dual exchange; \(E[d]\) perverse on a smooth \(d\)-fold.\\

M7 &
Fourier and curve engine. &
Delta/constant Fourier transforms, inversion shifts and twists.\\

M8 &
Direct-image weights and proper purity. &
No import from decomposition or final Weil files.\\

M9 &
IC purity, decomposition, and final assembly. &
All final declarations pass \(\#print\ \mathrm{axioms}\).\\
\bottomrule
\end{longtable}

\section{Import and axiom discipline}

Use four one-way import layers:
\[
\begin{tikzcd}[column sep=small]
\texttt{Generic}
 \arrow[r] &
\texttt{ProEtale}
 \arrow[r] &
\texttt{PerverseWeights}
 \arrow[r] &
\texttt{Weil}.
\end{tikzcd}
\tag{C.15}\label{eq:formalization-import-layers}
\]
No file in the first three layers may import the terminal Weil theorem.
Interface records used for parallel development must live under a visibly
temporary namespace and must not occur in the release import cone.

For every terminal theorem:

\begin{enumerate}
\item search the transitive imports for \(\texttt{sorry}\),
      \(\texttt{admit}\), and project-specific axioms;
\item run \(\#print\ \mathrm{axioms}\);
\item verify that all compactification, stratification, lattice, and finite
      coefficient-field choices have independence theorems;
\item test Frobenius and Tate conventions on the three basic examples above;
\item ensure that no definitional equality is being used where the
      mathematics supplies only a canonical isomorphism.
\end{enumerate}

This architecture makes the pro-etale and perverse infrastructure useful for
many other formalization projects: constructible sheaves with topological
coefficients, intersection cohomology, Fourier transforms, and decomposition
for proper maps all sit below the final Weil-specific layer.
